% 核心结论：去掉不一致的附加场与电荷正负示意，改用侧壁守恒与切向稳态
% 草图：[真实几何与节点 | 条件、推导、结论]；17cm印宽，最小字号8pt。
% 图契约：表面电荷调节法向场，使电流沿导线；图中的端口、几何、箭头和方程对应本题。
% 版面：沿用前三模块的几何图+推导分区；目标印宽17cm，简图至中等密度。
% 布局：左侧真实结构，右侧关键关系；真实电路边不作为架构流程箭头处理。
\documentclass[tikz,border=7pt]{standalone}
\usepackage{fontspec,xeCJK,amsmath,amssymb}
\IfFontExistsTF{Times New Roman}{\setmainfont{Times New Roman}}{\setmainfont{TeX Gyre Termes}}
\IfFontExistsTF{Heiti SC}{\setCJKmainfont[AutoFakeBold=1.5]{Heiti SC}}{\setCJKmainfont[AutoFakeBold=1.5]{FandolHei-Regular}}
\usetikzlibrary{arrows.meta,bending,calc,patterns,positioning,decorations.pathmorphing}
\definecolor{ink}{HTML}{243745}
\definecolor{blue}{HTML}{275DAB}
\definecolor{green}{HTML}{177C63}
\definecolor{red}{HTML}{BD4935}
\definecolor{orange}{HTML}{AA6C16}
\definecolor{purple}{HTML}{715393}
\tikzset{every node/.style={font=\fontsize{9.3}{12.5}\selectfont,text=ink},
arrow/.style={-{Stealth[length=5pt 1.5,width'=0pt 3]},shorten >=2pt,shorten <=1pt,line width=1pt},
arrow short/.style={-{Stealth[length=3pt 1.0,width'=0pt 2.5]},shorten >=1pt,shorten <=0pt,line width=.8pt},
axis/.style={arrow,draw=ink!60,line width=.6pt},
vec/.style={arrow,draw=blue},
wire/.style={draw=ink!75,line width=.8pt},
guide/.style={dashed,draw=ink!45,line width=.5pt},
card/.style={rounded corners=2pt,draw=blue!65,fill=blue!4,align=center,text width=6.5cm,inner xsep=9pt,inner ysep=7pt},
title/.style={font=\fontsize{10.5}{14}\selectfont\bfseries,text width=15cm,align=center},
label/.style={font=\fontsize{8.4}{11}\selectfont,align=center},
note/.style={label,text width=6.1cm}}
\begin{document}
\begin{tikzpicture}[x=1cm,y=1cm]
\path[use as bounding box] (-8,-3.5) rectangle(8,3.65);
\node[title] at (0,3.25) {\quad 表面电荷调节法向场，使电流沿导线};
\begin{scope}[shift={(-4.05,0)}]
\draw[ink!35,line width=8pt](-2.4,1.2)--(.7,1.2)--(.7,-1.5);\draw[vec,blue](-1.7,1.2)--(-.4,1.2);\draw[vec,green](.7,.7)--(.7,-.9);\draw[vec,red](1.9,.1)--(.8,.1)node[midway,above]{$\mathbf E_{\text{修正}}$};\node[label]at(-1.2,.45){稳态沿导线切向};\node[note]at(0,-2.8){局部修正场抵消原有法向分量；场的整体方向由两端电势决定};
\end{scope}
\node[card,text width=6.5cm,draw=blue!65,fill=blue!4]at(3.9,1.8){绝缘侧壁：$\partial_t\sigma_s=\mathbf J_{\rm in}\cdot\mathbf n$};
\node[card,text width=6.5cm,draw=green!65,fill=green!4]at(3.9,0){稳态：$\mathbf J\cdot\mathbf n=0$};
\node[card,text width=6.5cm,draw=orange!65,fill=orange!4]at(3.9,-1.8){$\mathbf J=\sigma\mathbf E$，总电场沿侧壁切向};

\end{tikzpicture}\end{document}